% GENERATED FILE. Edit canonical lessons, then run npm run generate:books.
% canonical-source-hash: fnv1a64:5370ffe17ad210ca
% canonical-lessons: ML-C483-sevanam, ML-C483-kotukuvala, ML-C483-sucitvam, ML-C483-menu, ML-C483-veyirrar

\chapter{Service and Meals}
\label{ch:ml-service-and-meals}
\hlchaptermodality{483}

\begin{chapteropening}
\textbf{By the end of this chapter:} \emph{I can say a service, a mosquito net, hygiene, a menu and a waiter.}

Complete the last lesson of chapter 483: I can say a service, a mosquito net, hygiene, a menu and a waiter.
\end{chapteropening}

\section[\texorpdfstring{\ml{സേവനം} (sēvanaṁ)}{സേവനം (sēvanaṁ)}]{\ml{സേവനം} (sēvanaṁ) — a service}

\label{lesson:ML-C483-sevanam}



\begin{quote}
Before the new one: say the Malayalam for to refund, then the Malayalam for an arrangement.
\end{quote}

\subsection*{You'll want to know: \ml{സേവനം}}
\textbf{\ml{സേവനം}} — \emph{sēvanaṁ} — \textquotedblleft{}a service\textquotedblright{}.

\begin{quote}
\ml{ആ} \ml{ഹോട്ടലിലെ} \ml{സേവനം} \ml{വളരെ} \ml{നല്ലതായിരുന്നു}.
\end{quote}

\begin{quote}
\emph{ā hōṭṭalile sēvanaṁ vaḷare nallatāyirunnu.}
\end{quote}

\begin{quote}
The service at that hotel was very good.
\end{quote}

A bus service, a bank service and public service are all \ml{സേവനം}.

\begin{cousinweb}
\textbf{\ml{സേവനം}} (\emph{sēvanaṁ}) is borrowed from Sanskrit \emph{sevana}, \textquotedblleft{}serving, attending on\textquotedblright{}, from the root \emph{sev}, to serve (Monier-Williams).
\end{cousinweb}

\begin{grammarlens}[title={What you've built}]
One more word for getting around.
\end{grammarlens}

\subsection*{Your turn}
\begin{itemize}
\raggedright
  \item \emph{Say it:} \emph{sēvanaṁ}
  \item \emph{Say it:} \emph{sēvanaṁ}, once more
  \item \emph{Say it:} \emph{ā hōṭṭalile sēvanaṁ vaḷare nallatāyirunnu.} — The service at that hotel was very good.
  \item \emph{Recall:} say the Malayalam for to refund, then the Malayalam for an arrangement, then say \emph{sēvanaṁ} again
\end{itemize}

\begin{tcolorbox}[breakable,colback=teal!4,colframe=teal!35!black,title={Before you move on}]
What does \ml{സേവനം} mean? (\textquotedblleft{}A service\textquotedblright{}.) Say it once more.
\end{tcolorbox}
\section[\texorpdfstring{\ml{കൊതുകുവല} (kotukuvala)}{കൊതുകുവല (kotukuvala)}]{\ml{കൊതുകുവല} (kotukuvala) — a mosquito net}

\label{lesson:ML-C483-kotukuvala}



\begin{quote}
Before the new one: say the Malayalam for an arrangement, then the Malayalam for a service.
\end{quote}

\subsection*{You'll want to know: \ml{കൊതുകുവല}}
\textbf{\ml{കൊതുകുവല}} — \emph{kotukuvala} — \textquotedblleft{}a mosquito net\textquotedblright{}.

\begin{quote}
\ml{രാത്രി} \ml{കുട്ടികൾ} \ml{കൊതുകുവലയ്ക്കുള്ളിൽ} \ml{ഉറങ്ങി}.
\end{quote}

\begin{quote}
\emph{rātri kuṭṭikaḷ kotukuvalaykkuḷḷil uṟaṅṅi.}
\end{quote}

\begin{quote}
At night the children slept inside a mosquito net.
\end{quote}

\ml{കൊതുക്}, a mosquito, plus \ml{വല}, a net. A fishing net is also a \ml{വല}.

\begin{cousinweb}
\textbf{\ml{കൊതുകുവല}} (\emph{kotukuvala}) is a native Dravidian compound of \ml{കൊതുക്}, a mosquito (DEDR 2086), and \ml{വല}, a net (DEDR 5288).
\end{cousinweb}

\begin{grammarlens}[title={What you've built}]
One more word for getting around.
\end{grammarlens}

\subsection*{Your turn}
\begin{itemize}
\raggedright
  \item \emph{Say it:} \emph{kotukuvala}
  \item \emph{Say it:} \emph{kotukuvala}, once more
  \item \emph{Say it:} \emph{rātri kuṭṭikaḷ kotukuvalaykkuḷḷil uṟaṅṅi.} — At night the children slept inside a mosquito net.
  \item \emph{Recall:} say the Malayalam for an arrangement, then the Malayalam for a service, then say \emph{kotukuvala} again
\end{itemize}

\begin{tcolorbox}[breakable,colback=teal!4,colframe=teal!35!black,title={Before you move on}]
What does \ml{കൊതുകുവല} mean? (\textquotedblleft{}A mosquito net\textquotedblright{}.) Say it once more.
\end{tcolorbox}
\section[\texorpdfstring{\ml{ശുചിത്വം} (śucitvaṁ)}{ശുചിത്വം (śucitvaṁ)}]{\ml{ശുചിത്വം} (śucitvaṁ) — hygiene; cleanliness}

\label{lesson:ML-C483-sucitvam}



\begin{quote}
Before the new one: say the Malayalam for a service, then the Malayalam for a mosquito net.
\end{quote}

\subsection*{You'll want to know: \ml{ശുചിത്വം}}
\textbf{\ml{ശുചിത്വം}} — \emph{śucitvaṁ} — \textquotedblleft{}hygiene; cleanliness\textquotedblright{}.

\begin{quote}
\ml{അടുക്കളയിൽ} \ml{ശുചിത്വം} \ml{വളരെ} \ml{പ്രധാനമാണ്}.
\end{quote}

\begin{quote}
\emph{aṭukkaḷayil śucitvaṁ vaḷare pradhānamāṇŭ.}
\end{quote}

\begin{quote}
Hygiene is very important in a kitchen.
\end{quote}

A more formal word than \ml{വൃത്തി}, cleanliness; it is common on notices.

\begin{cousinweb}
\textbf{\ml{ശുചിത്വം}} (\emph{śucitvaṁ}) is borrowed from Sanskrit \emph{śucitva}, \textquotedblleft{}clearness, purity\textquotedblright{} (Monier-Williams), from \emph{śuci}, clean, with the abstract ending \emph{-tva}.
\end{cousinweb}

\begin{grammarlens}[title={What you've built}]
One more word for getting around.
\end{grammarlens}

\subsection*{Your turn}
\begin{itemize}
\raggedright
  \item \emph{Say it:} \emph{śucitvaṁ}
  \item \emph{Say it:} \emph{śucitvaṁ}, once more
  \item \emph{Say it:} \emph{aṭukkaḷayil śucitvaṁ vaḷare pradhānamāṇŭ.} — Hygiene is very important in a kitchen.
  \item \emph{Recall:} say the Malayalam for a service, then the Malayalam for a mosquito net, then say \emph{śucitvaṁ} again
\end{itemize}

\begin{tcolorbox}[breakable,colback=teal!4,colframe=teal!35!black,title={Before you move on}]
What does \ml{ശുചിത്വം} mean? (\textquotedblleft{}Hygiene; cleanliness\textquotedblright{}.) Say it once more.
\end{tcolorbox}
\section[\texorpdfstring{\ml{മെനു} (menu)}{മെനു (menu)}]{\ml{മെനു} (menu) — a menu}

\label{lesson:ML-C483-menu}



\begin{quote}
Before the new one: say the Malayalam for a mosquito net, then the Malayalam for hygiene.
\end{quote}

\subsection*{You'll want to know: \ml{മെനു}}
\textbf{\ml{മെനു}} — \emph{menu} — \textquotedblleft{}a menu\textquotedblright{}.

\begin{quote}
\ml{ഹോട്ടലിലെ} \ml{മെനുവിൽ} \ml{മീൻകറി} \ml{ഉണ്ടായിരുന്നു}.
\end{quote}

\begin{quote}
\emph{hōṭṭalile menuvil mīnkaṟi uṇṭāyirunnu.}
\end{quote}

\begin{quote}
There was fish curry on the hotel menu.
\end{quote}

A loan ending in -\ml{ഉ} adds -\ml{വിൽ} for in: \ml{മെനുവിൽ}.

\begin{cousinweb}
\textbf{\ml{മെനു}} (\emph{menu}) was borrowed from English \emph{menu}, the list of dishes in a restaurant.
\end{cousinweb}

\begin{grammarlens}[title={What you've built}]
One more word for getting around.
\end{grammarlens}

\subsection*{Your turn}
\begin{itemize}
\raggedright
  \item \emph{Say it:} \emph{menu}
  \item \emph{Say it:} \emph{menu}, once more
  \item \emph{Say it:} \emph{hōṭṭalile menuvil mīnkaṟi uṇṭāyirunnu.} — There was fish curry on the hotel menu.
  \item \emph{Recall:} say the Malayalam for a mosquito net, then the Malayalam for hygiene, then say \emph{menu} again
\end{itemize}

\begin{tcolorbox}[breakable,colback=teal!4,colframe=teal!35!black,title={Before you move on}]
What does \ml{മെനു} mean? (\textquotedblleft{}A menu\textquotedblright{}.) Say it once more.
\end{tcolorbox}
\section[\texorpdfstring{\ml{വെയിറ്റർ} (veyiṟṟar)}{വെയിറ്റർ (veyiṟṟar)}]{\ml{വെയിറ്റർ} (veyiṟṟar) — a waiter}

\label{lesson:ML-C483-veyirrar}



\begin{quote}
Before the new one: say the Malayalam for hygiene, then the Malayalam for a menu.
\end{quote}

\subsection*{You'll want to know: \ml{വെയിറ്റർ}}
\textbf{\ml{വെയിറ്റർ}} — \emph{veyiṟṟar} — \textquotedblleft{}a waiter\textquotedblright{}.

\begin{quote}
\ml{വെയിറ്റർ} \ml{ഞങ്ങൾക്ക്} \ml{ചൂടുള്ള} \ml{ചായ} \ml{കൊണ്ടുവന്നു}.
\end{quote}

\begin{quote}
\emph{veyiṟṟar ñaṅṅaḷkkŭ cūṭuḷḷa cāya koṇṭuvannu.}
\end{quote}

\begin{quote}
The waiter brought us hot tea.
\end{quote}

In a small eating place people call out \ml{ചേട്ടാ} to him instead.

\begin{cousinweb}
\textbf{\ml{വെയിറ്റർ}} (\emph{veyiṟṟar}) was borrowed from English \emph{waiter}, a person who serves food at tables.
\end{cousinweb}

\begin{grammarlens}[title={What you've built}]
One more word for getting around.
\end{grammarlens}

\subsection*{Your turn}
\begin{itemize}
\raggedright
  \item \emph{Say it:} \emph{veyiṟṟar}
  \item \emph{Say it:} \emph{veyiṟṟar}, once more
  \item \emph{Say it:} \emph{veyiṟṟar ñaṅṅaḷkkŭ cūṭuḷḷa cāya koṇṭuvannu.} — The waiter brought us hot tea.
  \item \emph{Recall:} say the Malayalam for hygiene, then the Malayalam for a menu, then say \emph{veyiṟṟar} again
\end{itemize}

\begin{tcolorbox}[breakable,colback=teal!4,colframe=teal!35!black,title={Before you move on}]
What does \ml{വെയിറ്റർ} mean? (\textquotedblleft{}A waiter\textquotedblright{}.) Say it once more.
\end{tcolorbox}
